\chapter{مبرهنة بومبييري--فينوغرادوف والتوزيع المتوسطي للأعداد الأولية}

\section{نطاق الفصل وحالته}

\noindent\textbf{حالة الفصل:} \textenglish{DRAFT}. أُنشئ هذا المتن بعد
إغلاق بوابة ما قبل التأليف بحكم \textenglish{PASS-FOR-AUTHORING}. يثبت
الفصل مبرهنة بومبييري--فينوغرادوف في صيغة دالة فون مانغولت، ثم يستنتج
الصيغ الموافقة لـ\(\vartheta\) و\(\pi\) ونتيجة ``تقريبًا كل الترديدات''.

الانتقال المفاهيمي من الفصل الثاني عشر جوهري. فمبرهنة سيغل--فالفش تعطي
تقديرًا فرديًا موحدًا لكل ترديد
\[
q\le (\log x)^A,
\]
أما هنا فنسمح للترديدات بأن تبلغ تقريبًا \(x^{1/2}\)، لكن بعد أخذ متوسط
الأخطاء على \(q\). النتيجة المركزية هي أنه لكل \(A>0\)، وبانتظام عندما
\[
Q\le \frac{x^{1/2}}{(\log x)^{A+3}},
\]
لدينا
\[
\sum_{q\le Q}
\max_{(a,q)=1}
\sup_{2\le y\le x}
\left|
\psi(y;q,a)-\frac{y}{\varphi(q)}
\right|
\ll_A \frac{x}{(\log x)^A}.
\]

الثابت الضمني في المسار المعتمد غير فعال؛ فالبرهان يستعمل مبرهنة
سيغل--فالفش للموصلات الصغيرة. لا يثبت الفصل فرضية إليوت--هالبرستام، ولا
يتجاوز مستوى التوزيع \(1/2\) في الصيغة العامة غير الموزونة، ولا يعالج
الفترات القصيرة أو مبرهنة باربان--دافنبورت--هالبرستام.

الملفات الحاكمة للتأليف هي:
\begin{itemize}
  \item \texttt{research/literature-reviews/chapter-13-bombieri-vinogradov-evidence.md}.
  \item \texttt{research/literature-reviews/chapter-13-bombieri-vinogradov-proof-map.md}.
  \item \texttt{docs/CHAPTER\_13\_LARGE\_SIEVE\_MEAN\_VALUE\_AUDIT\_2026-07-21.md}.
  \item \texttt{docs/CHAPTER\_13\_VAUGHAN\_IDENTITY\_AUDIT\_2026-07-21.md}.
  \item \texttt{docs/CHAPTER\_13\_TYPE\_I\_TYPE\_II\_MEAN\_VALUE\_AUDIT\_2026-07-21.md}.
  \item \texttt{docs/CHAPTER\_13\_CONDUCTOR\_PRINCIPAL\_BV\_AUDIT\_2026-07-21.md}.
  \item \texttt{docs/CHAPTER\_13\_LOGIC\_AUDIT\_2026-07-21.md}.
  \item \texttt{docs/CHAPTER\_13\_REFERENCE\_VERIFICATION\_2026-07-21.md}.
\end{itemize}

\subsection{نواتج التعلم}

بعد إتمام الفصل ينبغي أن يستطيع القارئ:
\begin{enumerate}
  \item تفسير الفرق بين الانتظام الفردي والانتظام المتوسطي في الترديد.
  \item استعمال الغربال الكبير للشخصيات البدائية بوزنه الصحيح.
  \item إثبات هوية فوغان بتفاف ديريشليه.
  \item رد حدود هوية فوغان إلى مجاميع من النمطين الأول والثاني.
  \item إثبات متراجحة بوليا--فينوغرادوف من تحويل فورييه المنتهي.
  \item اختيار معاملي القطع \(U,V\) في المجالات المختلفة لـ\(Q\).
  \item فصل الشخصية الرئيسية ورد الشخصيات المستحثة إلى موصلاتها.
  \item تحديد الموضع الذي تدخل فيه عدم فعالية مبرهنة سيغل--فالفش.
  \item اشتقاق النسخ الموافقة لـ\(\vartheta\) و\(\pi\).
  \item تحويل تقدير متوسط إلى عبارة ``تقريبًا كل الترديدات''.
\end{enumerate}

\section{الأخطاء ومستوى التوزيع}

لكل \((a,q)=1\)، نضع
\[
E(y;q,a)
=
\psi(y;q,a)-\frac{y}{\varphi(q)},
\]
ونعرف
\[
E^*(x,q)
=
\max_{(a,q)=1}
\sup_{2\le y\le x}|E(y;q,a)|.
\]

\begin{definition}[مستوى توزيع متوسط]
نقول، بصورة غير رسمية في هذا الفصل، إن الأوليات تملك مستوى توزيع
\(\theta\) إذا أمكن ضبط
\[
\sum_{q\le x^{\theta-o(1)}}E^*(x,q)
\]
بادخار لوغاريتمي اعتباطي. تثبت مبرهنة بومبييري--فينوغرادوف المستوى
\(\theta=1/2\) حتى خسارة من قوة اللوغاريتم.
\end{definition}

لا تعني هذه العبارة أن كل ترديد مفرد حتى \(x^{1/2}\) يحقق تقديرًا من نوع
سيغل--فالفش. قد توجد ترديدات ذات خطأ أكبر، لكن مجموع الأخطاء على الترديدات
يبقى صغيرًا.

\section{حزمة الغربال الكبير}

\begin{theorem}[الغربال الكبير للشخصيات والنسخة الثنائية العظمى]
\label{thm:chapter13-large-sieve-package}
\resultid{ANT-THM-13-01}
\citedresult{\cite{Bombieri1965LargeSieve,MontgomeryVaughan2026}}
لتكن \(Q\ge1\).

أولًا، لكل معاملات مركبة \(a_n\):
\[
\sum_{q\le Q}\frac{q}{\varphi(q)}
\sum_{\chi\bmod q}^{*}
\left|
\sum_{M<n\le M+N}a_n\chi(n)
\right|^2
\le
(N+Q^2)
\sum_{M<n\le M+N}|a_n|^2.
\tag{13.1}
\]

ثانيًا، لكل متتاليتين مركبتين \(a_m,b_n\):
\begin{align*}
&\sum_{q\le Q}\frac{q}{\varphi(q)}
\sum_{\chi\bmod q}^{*}
\sup_Y
\left|
\sum_{\substack{m\le M,\ n\le N\\mn\le Y}}
 a_mb_n\chi(mn)
\right|\\
&\qquad\ll
(M+Q^2)^{1/2}(N+Q^2)^{1/2}
\left(\sum_{m\le M}|a_m|^2\right)^{1/2}
\left(\sum_{n\le N}|b_n|^2\right)^{1/2}
\log(2MN).
\tag{13.2}
\end{align*}
والنجمة تعني الجمع على الشخصيات البدائية.
\end{theorem}

الصيغة \((13.1)\) هي صورة الغربال الكبير التربيعية، أما \((13.2)\) فهي
حزمة مركبة تستخرج منها بكوشي وأداة قطع عظمى. لا ندعي في هذا الفصل برهانًا
داخليًا كاملًا للحزمة؛ بل نستعملها مدخلًا مقتبسًا مضبوط المنشأ.

\section{متراجحة بوليا--فينوغرادوف}

\begin{lemma}[بوليا--فينوغرادوف للشخصيات البدائية]
\label{lem:chapter13-polya-vinogradov}
\resultid{ANT-LEM-13-04}
\provedhere
إذا كانت \(\chi\) شخصية بدائية غير رئيسية modulo \(q>1\)، فإن لكل فترة
صحيحة
\[
M<n\le M+N
\]
لدينا
\[
\left|\sum_{M<n\le M+N}\chi(n)\right|
\ll \sqrt q\log(2q),
\tag{13.3}
\]
بثابت مطلق.
\end{lemma}

\begin{proof}
من تحويل فورييه المنتهي للشخصية البدائية وقيمة مجموع غاوس المطلقة في الفصل
السابع، نحصل -- مع مراعاة اصطلاح المرافق -- على صيغة عكسية من الشكل
\[
\chi(n)
=
\frac1{\tau(\overline\chi)}
\sum_{a=1}^{q}\overline\chi(a)e(an/q),
\qquad
|\tau(\overline\chi)|=\sqrt q.
\]
لذلك
\[
\left|\sum_{n\le N}\chi(n)\right|
\le
\frac1{\sqrt q}
\sum_{a=1}^{q-1}
\left|\sum_{n\le N}e(an/q)\right|.
\]
ويعطي تقدير المتسلسلة الهندسية
\[
\left|\sum_{n\le N}e(an/q)\right|
\ll
\min\left(N,\frac1{\|a/q\|}\right).
\]
وبتناظر \(a\) و\(q-a\):
\[
\sum_{a=1}^{q-1}
\min\left(N,\frac1{\|a/q\|}\right)
\ll
q\sum_{1\le a\le q/2}\frac1a
\ll q\log(2q).
\]
فتكون مجاميع البداية \(O(\sqrt q\log(2q))\)، ويعطي طرح مجموعي بدايتين
التقدير نفسه لأي فترة.
\end{proof}

\section{هوية فوغان}

لكل \(U,V\ge1\)، نضع
\[
\Lambda_U(n)=\Lambda(n)\mathbf1_{n\le U},
\qquad
\mu_V(n)=\mu(n)\mathbf1_{n\le V}.
\]

\begin{lemma}[هوية فوغان]
\label{lem:chapter13-vaughan-identity}
\resultid{ANT-LEM-13-01}
\provedhere
لكل \(n\ge1\):
\[
\Lambda(n)=c_1(n)+c_2(n)+c_3(n)+c_4(n),
\tag{13.4}
\]
حيث
\begin{align*}
c_1(n)
&=\Lambda(n)\mathbf1_{n\le U},\\
c_2(n)
&=-\sum_{\substack{rdk=n\\d\le U\\k\le V}}
\Lambda(d)\mu(k),\\
c_3(n)
&=\sum_{\substack{mk=n\\k\le V}}
\mu(k)\log m,\\
c_4(n)
&=\sum_{\substack{mk=n\\m>U\\k>V}}
\left(\sum_{\substack{d\mid m\\d>U}}\Lambda(d)\right)
\mu(k).
\end{align*}
\end{lemma}

\begin{proof}
على مستوى الدوال الحسابية والتفاف ديريشليه، تكون الحدود الأربعة
\[
\Lambda_U,
\qquad
-\mathbf1*\Lambda_U*\mu_V,
\qquad
\log*\mu_V,
\qquad
(\log-\mathbf1*\Lambda_U)*(\mu-\mu_V).
\]
بالتوسع والتجميع:
\begin{align*}
c_1+c_2+c_3+c_4
&=
\Lambda_U+\log*\mu-\mathbf1*\Lambda_U*\mu.
\end{align*}
ومن هويتي
\[
\log=\mathbf1*\Lambda,
\qquad
\mathbf1*\mu=\varepsilon,
\]
نحصل على
\[
\log*\mu=\Lambda,
\qquad
\mathbf1*\Lambda_U*\mu=\Lambda_U.
\]
إذن مجموع الحدود هو \(\Lambda\).

أما صيغة \(c_4\) المفصلة، فتنتج من
\[
\log m-
\sum_{\substack{d\mid m\\d\le U}}\Lambda(d)
=
\sum_{\substack{d\mid m\\d>U}}\Lambda(d).
\]
هذا العامل ينعدم عند \(m\le U\)، كما ينعدم \(\mu-\mu_V\) عند \(k\le V\)،
فتظهر شروط الدعم في \(c_4\) تلقائيًا.
\end{proof}

\section{الرد إلى النمطين الأول والثاني}

لتكن \(f\) دالة حسابية مركبة، ونضع
\[
S(Y;f)=\sum_{n\le Y}\Lambda(n)f(n).
\]

\begin{proposition}[تفكيك مجاميع فون مانغولت]
\label{prop:chapter13-type-decomposition}
\resultid{ANT-PROP-13-01}
\provedhere
تكتب \(S(Y;f)\) مجموع أربعة حدود، منها حد قصير، وحدان من النمط الأول، وحد
ثنائي من النمط الثاني:
\begin{align*}
S_1(Y)
&=\sum_{n\le\min(Y,U)}\Lambda(n)f(n),\\
S_2(Y)
&=\sum_{t\le UV}a(t)
\sum_{r\le Y/t}f(tr),
\qquad |a(t)|\le\log t,\\
S_3(Y)
&=\sum_{k\le V}\mu(k)
\sum_{m\le Y/k}f(km)\log m,\\
S_4(Y)
&=\sum_{\substack{mk\le Y\\m>U\\k>V}}
 b(m)\mu(k)f(mk),
\qquad |b(m)|\le\log m.
\end{align*}
\end{proposition}

\begin{proof}
نعوض هوية فوغان في تعريف \(S\)، ثم نضع
\[
a(t)=-\sum_{\substack{dk=t\\d\le U\\k\le V}}
\Lambda(d)\mu(k),
\qquad
b(m)=\sum_{\substack{d\mid m\\d>U}}\Lambda(d).
\]
ومن هوية فون مانغولت
\[
\sum_{d\mid t}\Lambda(d)=\log t
\]
نجد \(|a(t)|\le\log t\) و\(|b(m)|\le\log m\). ويكون \(S_2,S_3\) من
النمط الأول لأن أحد المتغيرين قصير، بينما يكون \(S_4\) ثنائيًا لأن
\(m>U\) و\(k>V\).
\end{proof}

\section{متوسط مجاميع الشخصيات}

لكل عائلة \(T(Y,\chi)\)، نضع
\[
\mathcal M(T)
=
\sum_{q\le Q}\frac{q}{\varphi(q)}
\sum_{\chi\bmod q}^{*}
\sup_{Y\le x}|T(Y,\chi)|.
\]

\subsection{تقدير النمط الأول}

\begin{lemma}[تقديرات Type I]
\label{lem:chapter13-type-one}
\resultid{ANT-LEM-13-02}
\provedhere
في تفكيك فوغان مع \(f(n)=\chi(n)\)، إذا قسمنا \(S_2=S_2'+S_2''\)
بحسب \(t\le U\) أو \(t>U\)، فإن
\[
\mathcal M(S_2')
\ll
\left(x+Q^{5/2}U\right)(\log(2Qx))^2,
\tag{13.5}
\]
و
\[
\mathcal M(S_3)
\ll
\left(x+Q^{5/2}V\right)(\log(2Qx))^2.
\tag{13.6}
\]
\end{lemma}

\begin{proof}
تعالج الشخصية الوحيدة modulo \(1\) تافهًا بمساهمة
\(O(x(\log 2x)^2)\). ولـ\(q>1\)، تعطي بوليا--فينوغرادوف
\[
\sup_z\left|\sum_{r\le z}\chi(rt)\right|
\le |\chi(t)|\sqrt q\log(2q).
\]
وباستعمال \(|a(t)|\le\log t\):
\[
\sup_{Y\le x}|S_2'(Y,\chi)|
\ll U\sqrt q\,(\log(2qU))^2.
\]
عدد الشخصيات البدائية لا يتجاوز \(\varphi(q)\)، ومع الوزن
\(q/\varphi(q)\) يظهر
\[
\sum_{q\le Q}q^{3/2}\ll Q^{5/2},
\]
فتنتج \((13.5)\).

أما \(S_3\)، فالجمع الجزئي يرد عامل \(\log m\) إلى مجاميع جزئية للشخصية
مع خسارة \(O(\log x)\). ثم تعطي بوليا--فينوغرادوف، بعد الجمع على
\(k\le V\)، الصيغة \((13.6)\).
\end{proof}

\subsection{تقدير النمط الثاني}

\begin{lemma}[تقديرات Type II]
\label{lem:chapter13-type-two}
\resultid{ANT-LEM-13-03}
\provedhere
لدينا
\[
\mathcal M(S_4)
\ll
\left(
x+QxU^{-1/2}+QxV^{-1/2}+Q^2x^{1/2}
\right)(\log 2x)^3,
\tag{13.7}
\]
كما أن الجزء \(S_2''\) يحقق
\[
\mathcal M(S_2'')
\ll
\left(
x+QxU^{-1/2}+Qx^{1/2}(UV)^{1/2}+Q^2x^{1/2}
\right)(\log 2x)^3.
\tag{13.8}
\]
\end{lemma}

\begin{proof}
نقسم المتغير الطويل إلى كتل ديادية \(M<m\le2M\). في كتلة من \(S_4\)،
تكون معياريات المعاملات
\[
\sum_m|b(m)|^2\ll M(\log 2x)^2,
\qquad
\sum_{k\le x/M}|\mu(k)|^2\le x/M.
\]
بتطبيق النسخة الثنائية العظمى \((13.2)\):
\[
\mathcal M(S_{4,M})
\ll
(M+Q^2)^{1/2}(x/M+Q^2)^{1/2}
 x^{1/2}(\log 2x)^2.
\]
وبتوسيع حاصل الضرب تحت الجذر، ثم جمع الكتل، تهيمن النهاية الدنيا \(U\)
على حدود \(M^{-1/2}\)، والنهاية العليا \(x/V\) على حدود \(M^{1/2}\)،
فتنتج \((13.7)\).

ولـ\(S_2''\)، نطبق الحساب نفسه مع المعاملين \(a(t)\) و\(1\)، وعلى المجال
\(U<t\le UV\). فتكون النهاية العليا \(UV\)، وينتج الحد
\(Qx^{1/2}(UV)^{1/2}\) في \((13.8)\).
\end{proof}

\begin{theorem}[مبرهنة القيمة المتوسطة للشخصيات البدائية]
\label{thm:chapter13-character-mean-value}
\resultid{ANT-THM-13-02}
\provedhere
لكل \(Q\ge1\) و\(x\ge2\):
\[
\sum_{q\le Q}\frac{q}{\varphi(q)}
\sum_{\chi\bmod q}^{*}
\sup_{Y\le x}|\psi(Y,\chi)|
\ll
\left(x+x^{5/6}Q+x^{1/2}Q^2\right)(\log 2x)^3.
\tag{13.9}
\]
\end{theorem}

\begin{proof}
من حد تشيبيشيف، يكون
\[
\mathcal M(S_1)\ll Q^2U,
\]
ويمتص هذا في \(Q^{5/2}U\). بجمع اللمتين السابقتين نحصل على
\begin{align*}
\mathcal M(\psi)
\ll
\Bigl(&x+QxU^{-1/2}+QxV^{-1/2}+Q^2x^{1/2}\\
&+Qx^{1/2}(UV)^{1/2}+Q^{5/2}U+Q^{5/2}V\Bigr)
(\log(2xUV))^3.
\tag{13.10}
\end{align*}

إذا \(Q\le x^{1/3}\)، نأخذ \(U=V=x^{1/3}\). وإذا
\(x^{1/3}\le Q\le x^{1/2}\)، نأخذ
\[
U=V=x^{2/3}/Q.
\]
في الحالتين تمتص كل حدود \((13.10)\) في الطرف الأيمن من \((13.9)\).

إذا \(Q>x^{1/2}\)، نطبق \((13.2)\) مباشرة مع
\[
M=1,
\qquad a_1=1,
\qquad b_n=\Lambda(n)\mathbf1_{n\le x}.
\]
ومن
\[
\sum_{n\le x}\Lambda(n)^2
\le(\log x)\psi(x)
\ll x\log(2x),
\]
نحصل على حد يغطّيه \(Q^2x^{1/2}(\log 2x)^3\). وهذا يكمل البرهان.
\end{proof}

\section{من الشخصيات إلى الفئات الحسابية}

لكل شخصية \(\chi\pmod q\)، نضع
\[
\psi'(Y,\chi)
=
\psi(Y,\chi)-\mathbf1_{\chi=\chi_0}Y.
\]
من تعامد الشخصيات:
\[
E(Y;q,a)
=
\frac1{\varphi(q)}
\sum_{\chi\bmod q}
\overline{\chi(a)}\psi'(Y,\chi),
\tag{13.11}
\]
ومن ثم
\[
E^*(x,q)
\le
\frac1{\varphi(q)}
\sum_{\chi\bmod q}
\sup_{Y\le x}|\psi'(Y,\chi)|.
\tag{13.12}
\]

إذا كانت \(\chi\pmod q\) مستحثة من شخصية بدائية \(\chi^*\pmod d\)، فإن
الفرق مدعوم على القوى الأولية المحذوفة، ولذلك
\[
|\psi'(Y,\chi)-\psi'(Y,\chi^*)|
\ll(\log(2qY))^2.
\tag{13.13}
\]

نحتاج أيضًا إلى اللمّة الأولية
\[
\sum_{m\le X}\frac1{\varphi(m)}\ll\log(2X).
\tag{13.14}
\]
فهي تنتج من الهوية
\[
\frac n{\varphi(n)}
=
\sum_{r\mid n}\frac{\mu^2(r)}{\varphi(r)}
\]
وتقارب
\[
\sum_{r\ge1}\frac{\mu^2(r)}{r\varphi(r)}.
\]
وبما أن \(\varphi(dm)\ge\varphi(d)\varphi(m)\)، نحصل على
\[
\sum_{\substack{q\le Q\\d\mid q}}
\frac1{\varphi(q)}
\ll
\frac1{\varphi(d)}\log\frac{2Q}{d}.
\tag{13.15}
\]

من \((13.12)\)--\((13.15)\):
\begin{align*}
\sum_{q\le Q}E^*(x,q)
&\ll
Q(\log(2Qx))^2\\
&\quad+
\sum_{d\le Q}
\frac{\log(2Q/d)}{\varphi(d)}
\sum_{\chi\bmod d}^{*}
\sup_{Y\le x}|\psi'(Y,\chi)|.
\tag{13.16}
\end{align*}

\section{مبرهنة بومبييري--فينوغرادوف}

\begin{theorem}[بومبييري--فينوغرادوف]
\label{thm:chapter13-bombieri-vinogradov}
\resultid{ANT-THM-13-03}
\provedhere
لكل \(A>0\)، وبانتظام عندما
\[
Q\le\frac{x^{1/2}}{(\log x)^{A+3}},
\]
لدينا
\[
\sum_{q\le Q}
\max_{(a,q)=1}
\sup_{2\le Y\le x}
\left|
\psi(Y;q,a)-\frac{Y}{\varphi(q)}
\right|
\ll_A
\frac{x}{(\log x)^A}.
\tag{13.17}
\]
الثابت الضمني غير فعال في المسار المعتمد.
\end{theorem}

\begin{proof}
نضع
\[
L=\log x,
\qquad
D=L^{A+4}.
\]
نفصل مجموع \((13.16)\) عند \(d=D\).

\paragraph{الموصلات الصغيرة.}
إذا \(d\le D\)، فإن مبرهنة سيغل--فالفش من الفصل الثاني عشر، مع أس
كبير بما يكفي، تعطي لكل شخصية بدائية غير رئيسية
\[
\sup_{Y\le x}|\psi(Y,\chi)|
\ll_A xe^{-c_A\sqrt L}.
\tag{13.18}
\]
وللشخصية ذات الموصل \(1\)، تعطي مبرهنة الأعداد الأولية الفعالة التقدير
نفسه لـ\(|\psi(Y)-Y|\). لتوحيد القيمة العظمى، نطبق سيغل--فالفش عندما
\(Y\ge e^{\sqrt L}\)، ونعالج \(Y<e^{\sqrt L}\) تافهًا. لذلك تكون مساهمة
\(d\le D\)
\[
\ll_A D\log(2Q)xe^{-c_A\sqrt L}
\ll_A xL^{-A}.
\tag{13.19}
\]
هذا هو موضع عدم الفعالية.

\paragraph{الموصلات الكبيرة.}
نقسم \(D<d\le Q\) إلى كتل ديادية \(U<d\le2U\). في هذه الكتل تكون
الشخصيات البدائية غير رئيسية، ومن
\[
\frac1{\varphi(d)}
\le
\frac1U\frac d{\varphi(d)}
\]
ومبرهنة القيمة المتوسطة نحصل على
\begin{align*}
&\sum_{U<d\le2U}
\frac{\log(2Q/d)}{\varphi(d)}
\sum_{\chi\bmod d}^{*}
\sup_{Y\le x}|\psi(Y,\chi)|\\
&\qquad\ll
\left(
\frac xU+x^{5/6}+x^{1/2}U
\right)L^3\log(4Q/U).
\end{align*}
بعد جمع الكتل تكون المساهمة
\[
\ll
xD^{-1}L^4+x^{5/6}L^5+x^{1/2}QL^3.
\tag{13.20}
\]

نأخذ
\[
Q_0=x^{1/2}L^{-(A+3)}.
\]
عندئذ
\[
xD^{-1}L^4=xL^{-A},
\qquad
x^{1/2}Q_0L^3=xL^{-A},
\]
كما أن \(x^{5/6}L^5\ll_A xL^{-A}\). ويمتص الخطأ المحلي الأول في
\((13.16)\) أيضًا. فتنتج \((13.17)\) عند \(Q_0\)، ثم تمتد إلى كل
\(Q\le Q_0\) برتابة الطرف الأيسر.
\end{proof}

\section{نتائج تابعة}

\begin{corollary}[النسخة الموافقة لـ\(\vartheta\)]
\label{cor:chapter13-theta}
\resultid{ANT-COR-13-01}
\provedhere
لكل \(A>0\)، إذا
\[
Q\le\frac{x^{1/2}}{(\log x)^{A+3}},
\]
فإن
\[
\sum_{q\le Q}
\max_{(a,q)=1}
\sup_{2\le Y\le x}
\left|
\vartheta(Y;q,a)-\frac{Y}{\varphi(q)}
\right|
\ll_A\frac{x}{(\log x)^A}.
\tag{13.21}
\]
\end{corollary}

\begin{proof}
من ضبط القوى الأولية العليا في الفصل التاسع:
\[
|\psi(Y;q,a)-\vartheta(Y;q,a)|
\le\psi(Y)-\vartheta(Y)
\ll Y^{1/2}\log(2Y).
\]
وبجمع هذا الخطأ على \(q\le Q\)، نحصل على
\[
Qx^{1/2}\log(2x)
\le x(\log x)^{-(A+2)},
\]
فيمتص في \((13.17)\).
\end{proof}

\begin{corollary}[النسخة الموافقة لـ\(\pi\)]
\label{cor:chapter13-pi}
\resultid{ANT-COR-13-02}
\provedhere
لكل \(A>0\)، إذا
\[
Q\le\frac{x^{1/2}}{(\log x)^{A+4}},
\]
فإن
\[
\sum_{q\le Q}
\max_{(a,q)=1}
\sup_{2\le Y\le x}
\left|
\pi(Y;q,a)-\frac{\operatorname{Li}(Y)}{\varphi(q)}
\right|
\ll_A\frac{x}{(\log x)^A}.
\tag{13.22}
\]
\end{corollary}

\begin{proof}
بالتجميع الجزئي:
\[
\pi(Y;q,a)
=
\frac{\vartheta(Y;q,a)}{\log Y}
+
\int_2^Y
\frac{\vartheta(t;q,a)}{t(\log t)^2}\,dt.
\]
بعد طرح \(\operatorname{Li}(Y)/\varphi(q)\)، يضبط الخطأ بأكبر خطأ في
\(\vartheta\) مضافًا إليه حد مطلق. نطبق النتيجة السابقة بأس الادخار
\(A+1\)، ولذلك يصبح مجال الترديد \(A+4\). أما مجموع الحدود المطلقة على
\(q\le Q\) فيمتص في الطرف المطلوب.
\end{proof}

\begin{corollary}[تقريبًا كل الترديدات]
\label{cor:chapter13-almost-all-moduli}
\resultid{ANT-COR-13-03}
\provedhere
لتكن \(B,C>0\)، وافترض
\[
Q\le\frac{x^{1/2}}{(\log x)^{B+C+3}}.
\]
عندئذ، جميع الترديدات \(q\le Q\) عدا
\[
O_{B,C}\left(\frac{Q}{(\log x)^C}\right)
\]
منها تحقق
\[
E^*(x,q)
\le
\frac{x}{Q(\log x)^B}.
\tag{13.23}
\]
\end{corollary}

\begin{proof}
نسمي \(q\) سيئًا إذا خالف \((13.23)\)، وليكن عدد الترديدات السيئة \(R\).
بتطبيق مبرهنة بومبييري--فينوغرادوف بالأس \(B+C\):
\[
R\frac{x}{Q(\log x)^B}
\le
\sum_{q\le Q}E^*(x,q)
\ll_{B,C}
\frac{x}{(\log x)^{B+C}}.
\]
ومن ثم \(R\ll Q(\log x)^{-C}\).
\end{proof}

\section{مقارنة النتائج الثلاث}

\begin{center}
\begin{tabular}{|c|c|c|}
\hline
النتيجة & نوع الانتظام & مجال الترديد \\
\hline
الفصل العاشر & فردي، \(q\) ثابت & ثابت \\
\hline
سيغل--فالفش & فردي موحد & \(q\le(\log x)^A\) \\
\hline
بومبييري--فينوغرادوف & متوسط على \(q\) & \(q\le x^{1/2}(\log x)^{-A-3}\) \\
\hline
إليوت--هالبرستام & متوسط تخميني & \(q\le x^{1-\varepsilon}\) \\
\hline
\end{tabular}
\end{center}

مستوى \(1/2\) ليس حدًا منطقيًا مطلقًا لكل مسألة؛ توجد نتائج تتجاوزه بأوزان
خاصة أو بقيود بنيوية على الترديدات. لكن هذه النتائج لا تثبت فرضية
إليوت--هالبرستام العامة، ولذلك تبقى خارج النواة البرهانية لهذا الفصل.

\section{مثال تفسيري}

إذا أخذنا \(A=2\)، تسمح المبرهنة بالترديدات حتى
\[
Q\le\frac{x^{1/2}}{(\log x)^5},
\]
وتعطي
\[
\sum_{q\le Q}E^*(x,q)
\ll\frac{x}{(\log x)^2}.
\]
لا نستنتج من ذلك أن كل \(E^*(x,q)\) أصغر من
\(x/(Q(\log x)^2)\). لكن نتيجة تقريبًا كل الترديدات تبين أن عدد الترديدات
ذات الخطأ الأكبر من عتبة مختارة صغير نسبيًا، بعد تخصيص ادخار لوغاريتمي
إضافي.

\section{حدود الادعاء}

\begin{itemize}
  \item حزمة الغربال الكبير \((13.1)\)--\((13.2)\) مقتبسة، لا مثبتة من
  الصفر في هذا الفصل.
  \item هوية فوغان وPólya--Vinogradov وتقديرا Type I وType II ومبرهنة
  القيمة المتوسطة مثبتة داخل الفصل اعتمادًا على الحزمة المقتبسة.
  \item الثوابت النهائية غير فعالة بسبب استعمال سيغل--فالفش للموصلات
  الصغيرة.
  \item لا يثبت الفصل Elliott--Halberstam أو Barban--Davenport--Halberstam.
  \item لا يعالج الفصل الفترات القصيرة أو تطبيقات الفجوات المحدودة أو
  غولدباخ؛ هذه موضوعات لاحقة.
  \item لا تعني حالة \textenglish{DRAFT} أن الفصل اجتاز التحقق المرجعي
  للمتن أو تدقيق ما بعد التأليف أو المراجعة المستقلة.
\end{itemize}

\section{أسئلة تفسير وتمارين}

\begin{enumerate}
  \item اشرح لماذا لا تناقض مبرهنة بومبييري--فينوغرادوف احتمال وجود ترديد
  مفرد ذي خطأ كبير نسبيًا.
  \item تحقق مباشرة من هوية فوغان عند \(n=p\)، وعند \(n=p^2\)، لعدة علاقات
  بين \(p,U,V\).
  \item أثبت الهوية
  \[
  \frac n{\varphi(n)}
  =
  \sum_{d\mid n}\frac{\mu^2(d)}{\varphi(d)}.
  \]
  \item أكمل تفاصيل الجمع الجزئي في تقدير \(S_3\).
  \item تحقق من كل متراجحات اختيار \(U,V\) في المجال
  \(x^{1/3}\le Q\le x^{1/2}\).
  \item بيّن أن
  \[
  \sum_{n\le x}\Lambda(n)^2\ll x\log x
  \]
  باستعمال حد تشيبيشيف.
  \item استخرج من \((13.17)\) عبارة ``تقريبًا كل الترديدات'' بعتبة مختلفة
  عن \((13.23)\).
  \item اشرح بدقة لماذا يؤدي استعمال سيغل--فالفش في الموصلات الصغيرة إلى
  عدم فعالية الثابت النهائي.
  \item قارن منطقيًا بين مبرهنة بومبييري--فينوغرادوف وفرضية
  إليوت--هالبرستام، من دون الخلط بين نتيجة موزونة تتجاوز \(1/2\) والفرضية
  العامة.
\end{enumerate}

\section{خلاصة الفصل}

بدأ الفصل من غربال كبير مقتبس للشخصيات البدائية، وأثبت داخليًا هوية فوغان
ومتراجحة بوليا--فينوغرادوف، ثم ضبط مجاميع النمطين الأول والثاني، فحصل على
مبرهنة قيمة متوسطة لمجاميع فون مانغولت الملتوية. بعد رد الشخصيات المستحثة
إلى موصلاتها وفصل الشخصية الرئيسية والموصلات الصغيرة، نتجت مبرهنة
بومبييري--فينوغرادوف عند مستوى توزيع \(1/2\) حتى خسارة لوغاريتمية. وأخيرًا
انتقل البرهان إلى \(\vartheta\) و\(\pi\) واستخرج نتيجة تقريبًا كل
الترديدات، مع بقاء عدم الفعالية وحدود النطاق معلنة بوضوح.
